% ABSTRACT-START
\begin{abstract}
A recognition prefix, or secret handshake, lets a strategy in the iterated
prisoner's dilemma identify a partner before committing to cooperation: it plays
a fixed word, then cooperates only if the opponent has echoed that word. The
handshake idea and the structural invariants of binary words are each classical,
yet the prefix space itself has not been enumerated and profiled as an object.
We enumerate the deterministic recognition-prefix strategies and apply a joint
structural and evolutionary battery: minimal-automaton size, aperiodic
autocorrelation, self-overlap, sequence complexity, and Moran-process fixation
under three post-recognition regimes. Profiling reproduces the complete Barker
spectrum, with peak-sidelobe-one prefixes occurring at exactly
$k=2,3,4,5,7,11,13$ and at no other length up to $13$, a spectrum nowhere
encoded in the tool. The minimal Moore machine of a length-$k$ prefix takes
every value in the band $[k+2,\,2k+1]$ for all $k\le 11$. A single cooperative
bit decides whether a handshake can invade a population of defectors: the
all-defect prefix sits below neutral drift, while one $C$ bit raises fixation
above neutral. Mimic-resistance is structure-independent under cooperate-forever
recognition, where a forger always invades, but vigilant grim and tit-for-tat
regimes reverse this for every handshake across all $36$ tested parameter cells.
The enumeration tool and all data are released open source for reproduction and
extension.
\end{abstract}
% ABSTRACT-END

\section{Introduction}
\label{sec:intro}

Cooperation in the iterated prisoner's dilemma (IPD) is fragile: a population of
cooperators is open to invasion by defectors who take the tempting payoff and
give nothing back \citep{axelrod1984}. One escape is recognition. If a strategy
can identify a like-minded partner before it commits to cooperation, it can
extend cooperation to kin and withhold it from everyone else. The classical
device for this is the secret handshake \citep{robson1990}: a strategy plays a
fixed sequence of moves, watches whether the opponent plays the same sequence,
and cooperates only on a match. The handshake is a recognition prefix, a word
$p\in\{C,D\}^k$ that gates the strategy's later behaviour.

Recognition prefixes are not a curiosity. They emerge spontaneously when finite-
state strategies are evolved in the Moran process: handshake-like prefixes
appear in the machines that survive, and they sharpen the line between strategies
that invade defectors and strategies that do not \citep{knight2018}. The prefix
is also a natural unit of strategy complexity. A strategy realised by a finite
automaton has a minimal-state form, and the number of states is the standard
measure of how complicated the strategy is \citep{rubinstein1986,abreu1988}; a
handshake is among the simplest constructions that still implements conditional
cooperation, and low-complexity structure is where direct reciprocity tends to
settle \citep{press2012,hilbe2018}.

The handshake concept and the mathematics it touches are both well established,
and they have been developed apart from each other. The minimal automaton of a
strategy and its computation by Myhill--Nerode equivalence and partition
refinement are textbook \citep{nerode1958,hopcroft1971,rubinstein1986}. The
autocorrelation properties of binary words, and the lengths at which a word has a
peak sidelobe of one, are the subject of the Barker-code literature
\citep{barker1953}. The self-overlap structure of words, including which words
have no proper border, has a combinatorial theory of its own
\citep{guibas1981}. What has not been done is to treat the recognition-prefix
space as a single object: to enumerate it exhaustively and to profile each prefix
with all of these invariants at once, together with how the gated strategy fares
evolutionarily. We do that here.

We are precise about the novelty. The handshake mechanism is from
\citet{robson1990}; the emergence of recognition in evolved machines and the
invasion-of-defectors test are from \citet{knight2018}; every structural
invariant we compute is a known primitive. The contribution is the synthesis:
the exhaustive enumeration of the prefix space and its profiling under a joint
structural and evolutionary battery, which lets us state where structure lives in
the space, how it scales, and when it matters for cooperation. We claim the
catalogue and what it reveals, not the concept or the primitives.

\begin{itemize}
  \item \textbf{An exhaustive enumeration with a joint invariant battery.} We
    enumerate all deterministic recognition-prefix strategies and compute, for
    each, a structural profile (minimal-automaton size, aperiodic
    autocorrelation, self-overlap and borderedness, and sequence complexity) and
    an evolutionary profile (Moran-process fixation under three post-recognition
    regimes). The enumeration is deterministic and the fixation probabilities are
    closed form, so the catalogue is exactly reproducible
    (\cref{sec:method}).
  \item \textbf{A reproducible family taxonomy.} We assign each prefix to one of
    six rule-based families by priority predicates over its invariants, and we
    cross-check the assignment against unsupervised clustering of the invariant
    vectors (\cref{tab:families}, \cref{sec:results}).
  \item \textbf{Findings.} Profiling recovers the complete Barker spectrum as a
    validation of the autocorrelation pipeline (\cref{tab:barker}); the minimal
    automaton occupies an exact band $[k+2,\,2k+1]$ for every enumerated length
    $k\le 11$ (\cref{tab:autoband}); a
    single cooperative bit is the threshold for a handshake to invade a defector
    population (\cref{tab:invade}); and mimic-resistance is structure-independent
    under cooperate-forever recognition but is restored for every handshake by the
    vigilant grim and tit-for-tat regimes (\cref{tab:regime},
    \cref{tab:sensitivity}).
\end{itemize}

\section{Preliminaries}
\label{sec:prelim}

\paragraph{Stage game and payoffs.}
Each round both players choose an action from $\{C,D\}$, where $C$ is cooperate
and $D$ is defect. The payoffs are $R=3$ for mutual cooperation, $P=1$ for mutual
defection, $T=5$ to a defector facing a cooperator, and $S=0$ to the exploited
cooperator. These satisfy $T>R>P>S$ and $2R>T+S$, the standard prisoner's-dilemma
ordering \citep{axelrod1984}. A match runs for $n$ rounds ($n=200$ by default),
and a player's score is the mean payoff per round.

\paragraph{Recognition-prefix strategies.}
A handshake is a recognition prefix $p\in\{C,D\}^k$ of length $k\ge 1$. The gated
strategy $H_p$ with post-recognition regime $\rho$ behaves as follows. For the
first $k$ rounds it plays $p$ regardless of what the opponent does. After round
$k$ it compares the opponent's first $k$ actions with $p$. If they are equal, the
opponent is recognised and $H_p$ enters regime $\rho$; otherwise $H_p$ defects on
every remaining round. The first $k$ rounds are the recognition phase, and the
prefix doubles as the signal $H_p$ emits and the password it requires.

\paragraph{Post-recognition regimes.}
We study three regimes $\rho\in\{\textsc{coop},\textsc{grim},\textsc{tft}\}$.
Under \textsc{coop} the strategy cooperates on every round after recognition; this
is the classical secret handshake of \citet{robson1990}, which trusts a
recognised partner unconditionally. Under \textsc{grim} the strategy cooperates
with a recognised partner until that partner defects once, then defects on every
remaining round. Under \textsc{tft} the strategy cooperates on the first post-
recognition round and thereafter copies the partner's previous action. The
regimes differ only in how a recognised partner is treated; recognition itself is
identical.

\paragraph{The mimic.}
A mimic of $p$ is the cheat that forges the signal: it plays $p$ through the
recognition phase, so $H_p$ recognises it, and then defects on every remaining
round. The mimic is the natural adversary for the handshake, and whether it can
invade a population of handshakers measures how much the recognition gate is
worth.

\paragraph{Minimal Moore machine.}
We realise $H_p$ as a Moore machine, a finite-state automaton that emits an
action in each state and transitions on the opponent's action. Two states are
equivalent when they induce the same behaviour against every continuation, the
Myhill--Nerode relation \citep{nerode1958}; collapsing equivalent states gives
the unique minimal machine, which we compute by Hopcroft partition refinement
\citep{hopcroft1971}. We write $m(p)$ for the number of states of this minimal
machine. It is the Rubinstein measure of strategy complexity
\citep{rubinstein1986}, the size of the smallest automaton that plays $H_p$.

\paragraph{Aperiodic autocorrelation.}
Map $p$ to its $\pm1$ image $a\in\{-1,+1\}^N$ with $N=k$ by sending $C$ to $+1$
and $D$ to $-1$. The aperiodic autocorrelation at shift $v$ is
$c_v=\sum_j a_j a_{j+v}$, summed over the overlapping positions. The sidelobes
are the $c_v$ for $v\neq 0$, the peak sidelobe is $\max_{v\neq 0}|c_v|$, and the
merit factor is $N^2/(2\sum_{v\neq 0} c_v^2)$. A word is a Barker code when its
peak sidelobe is at most $1$ \citep{barker1953}, the strongest possible
flatness, and such codes exist only at a short list of lengths
\citep{turyn1961}.

\paragraph{Self-overlap and bordered words.}
A border of $p$ is a proper prefix that is also a suffix. A word is unbordered,
or bifix-free, when it has no proper border, so it cannot align with a shifted
copy of itself except at full overlap. Borders and the correlation polynomial
that records them are the Guibas--Odlyzko theory of string overlaps
\citep{guibas1981}; the number of unbordered binary words of each length is the
sequence \citep{oeisA003000}.

\paragraph{Sequence complexity.}
We summarise the incompressibility of $p$ by its LZ76 production complexity, the
number of distinct factors a left-to-right parse produces \citep{lempel1976}, and
by its linear complexity, the length of the shortest linear-feedback shift
register that generates $p$, computed by the Berlekamp--Massey algorithm
\citep{massey1969}. Together these separate structured prefixes from
near-incompressible ones.

\paragraph{Moran process.}
We place $H_p$ in a finite population of size $N$ ($N=20$ by default) evolving
under the Moran process \citep{moran1958}: each step one individual reproduces
with probability proportional to fitness and one is removed at random. Fitness is
mapped from match score with selection intensity $w$ ($w=0.1$ by default). The
fixation probability of a single mutant in a resident population follows the
closed-form constant-selection formula \citep{nowak2004}, which we evaluate
exactly with no simulation. The neutral-drift baseline, the fixation probability
of a selectively neutral mutant, is $1/N$.
